% Generated draft: requires mathematical and source review.
\documentclass[11pt]{article}
\usepackage[margin=25mm]{geometry}
\usepackage{fontspec}
\setmainfont{Latin Modern Roman}
\usepackage{amsmath,amssymb,mathtools,unicode-math}
\setmathfont{Latin Modern Math}
\usepackage{enumitem,needspace,xurl,hyperref}
\setlength{\parindent}{0pt}
\setlength{\parskip}{5pt}
\setlength{\emergencystretch}{3em}
\newcommand{\N}{\mathbb{N}}
\newcommand{\Z}{\mathbb{Z}}
\newcommand{\Q}{\mathbb{Q}}
\newcommand{\R}{\mathbb{R}}
\newcommand{\C}{\mathbb{C}}
\newcommand{\F}{\mathbb{F}}
\newcommand{\A}{\mathbb{A}}
\newcommand{\PP}{\mathbb{P}}
\newcommand{\E}{\mathbb{E}}
\newcommand{\eps}{\varepsilon}
\newcommand{\dd}{\,\mathrm{d}}
\begin{document}
\section*{UnsolvedMath ID: 30003292}
\textbf{Source problem code:} catalogue entry
\subsection{Definitions and mathematical statement}
A group \(\Gamma\) is \emph{coherent} if every finitely generated subgroup of \(\Gamma\) is finitely presented. A finite \(K(G,1)\) means a connected CW complex with finitely many cells, fundamental group isomorphic to \(G\), and contractible universal cover.

For clarity, the following describes the usual integral interpretation of the homological terminology. A group \(H\) has type \(\mathrm{FP}_m\), for an integer \(m\geq 0\), if the trivial \(\Z H\)-module \(\Z\) admits a projective resolution whose modules in degrees \(0,\ldots,m\) are finitely generated. Type \(\mathrm{FP}\) means that it admits a finite-length resolution by finitely generated projective \(\Z H\)-modules. This is distinguished from type \(\mathrm{FP}_{\infty}\), which requires finite generation in every degree without requiring finite length. An integral Poincar\'e duality group of dimension \(n\), denoted \(\mathrm{PD}^n\), is a group \(P\) of type \(\mathrm{FP}\) such that
\[
H^i(P;\Z P)\cong
\begin{cases}
0,&i\neq n,\\
\Z,&i=n,
\end{cases}
\]
as abelian groups; the induced action on the degree-\(n\) dualizing module may be nontrivial. The source does not specify coefficients, orientability, or the precise meaning of its unindexed \(\mathrm{FP}\); these conventions therefore require confirmation against the report.

Consider a surface bundle
\[
F\longrightarrow E\longrightarrow B,
\qquad \pi=\pi_1(E,e),
\]
where \(B\) is a closed surface and \(e\in E\) is a chosen basepoint. Here closed means compact without boundary. The supplied record concerns surface bundles, but its wording does not separately specify whether \(F\) is closed, whether the surfaces are connected or orientable, or whether there are genus restrictions. The intended range of bundles must consequently be checked; none of these restrictions is silently imposed here.

The supplied three-part question is:
\begin{enumerate}
\item For every bundle in the intended class, is \(\pi\) coherent? Equivalently, does
\[
\forall H\leq\pi,\qquad H\text{ finitely generated}\ \Longrightarrow\ H\text{ finitely presented}
\]
hold?
\item For every such bundle and every finitely presented subgroup \(G\leq\pi\), does \(G\) admit a finite \(K(G,1)\)?
\item More generally, for every dimension \(n\), every \(\mathrm{PD}^n\) group \(P\), and every subgroup \(H\leq P\), does
\[
H\text{ of type }\mathrm{FP}_{\lfloor n/2\rfloor}
\quad\Longrightarrow\quad
H\text{ of type }\mathrm{FP}
\]
hold? The subscript notation here represents the source's notation \(\mathrm{FP}^{\lfloor n/2\rfloor}\). Under the integral interpretation above, \(n\) is a nonnegative integer; the source states no narrower dimension range.
\end{enumerate}
These are the recorded questions, subject to the unresolved scope and terminology qualifications above, and are not assertions that the implications hold.
\subsection{Short English statement}
The problem asks whether fundamental groups of the intended surface bundles are coherent and whether all their finitely presented subgroups have finite classifying spaces. It also asks whether a subgroup of a Poincar\'e duality group with the stated partial homological finiteness property must have type \(\mathrm{FP}\). The permitted surfaces and the homological conventions need confirmation. The supplied status is partially solved; this draft does not independently verify that assessment.
\subsection{Sources}
\begin{itemize}
\item Local UnsolvedMath registry, permanent ID \(30003292\), title: \emph{Coherence and Finiteness Properties of Surface-Bundle Groups}; exact source problem code: \texttt{catalogue entry}. The record belongs to the Oberwolfach Reports Open Problems collection. No individual proposer is supplied; the recorded proposed year is 2016.
\item Supplied source citation: \emph{Surface Bundles (2017)}, Oberwolfach Reports, DOI 10.4171/owr/2016/56, pp.\ 3149--3195. \url{https://doi.org/10.4171/owr/2016/56}. The background describes these reports as documenting mathematical workshops and their open-problem discussions.
\item Supplied literature source, labeled \emph{Oberwolfach Report 56 (2017)}: \url{https://doi.org/10.4171/owr/2016/56}.
\item Supplied literature source, labeled \emph{Baumslag and Roseblade subgroup theorem}: \url{https://people.maths.ox.ac.uk/~bridson/papers/BMiller07/BMiller08.pdf}. The record supplies no full bibliographic citation or theorem location for this link.
\item Literature assessment dated exactly \texttt{2026-08-21}: the supplied curation labels the problem partially solved and reports that coherence and finiteness properties are established in several cases, that non-coherence is forced in others, and that the full three-part question for arbitrary surface-by-surface groups remains open. This is supplied dated curation derived from the OWR material, not an independently verified literature assessment; its time-sensitive status and cited claims require independent checking.
\item The registry labels the statement \texttt{corrected\_verified} and reports a manual semantic, typographic, and LaTeX review of the clean statement. These are supplied editorial claims, not mathematical verification performed for this draft.
\end{itemize}
\end{document}
